\section{Boundary maps and Hall numbers}\label{sec:boundary}
All modules in the remaining sections are right modules. Products of paths are written in walking order: if $a:x\to y$ and $b:y\to z$, then $ab:x\to z$. Composition of linear maps retains the convention $gf=g\circ f$. At a quiver vertex $v$, write $e_v$ for the primitive idempotent and $S_v$ for the corresponding simple module. For a module $V$, put $\rad V=V\rad A$; its socle $\soc V$ is its largest semisimple submodule. We write $D_kV=\Hom_k(V,k)$.

We first give a counting argument for two algebras with a common quotient. Suppose $A_0,A_1$ have ideals $J_0,J_1$ and specified identifications $A_i/J_i=C$. Suppose corresponding modules $E_i$ have identified $C$-modules
\[
 S=E_iJ_i,\qquad R=\{v\in E_i:vJ_i=0\},\qquad T=E_i/S,
\]
where $S\subseteq R$. Let $h_i:R\to T$ be the induced maps, with kernel $S$. Assume their common image is a submodule $H\subseteq T$ and
\begin{equation}\label{eq:boundary-shear}
 h_1=\phi h_0,
\end{equation}
where $\phi\in\Aut_C(H)$ preserves every submodule of $H$. For $C$-modules $V,W$, let $\mathcal P_C(V,W)$ be the subspace of $\Hom_C(V,W)$ consisting of maps factoring through finite-dimensional projective $C$-modules.

\begin{lemma}\label{lem:graph-comparison}
Suppose that, for every $\Omega\subseteq H$, every $C$-module $N'$ and every $g:\Omega\to N'$, there is $\psi\in\Aut_C(N')$ such that
\begin{equation}\label{eq:stable-compensation}
 \psi g-g\phi|_\Omega\in\mathcal P_C(\Omega,N').
\end{equation}
Then for all $C$-modules $X,Y,N$,
\[
 F^{E_0\oplus N}_{X,Y}(A_0)=F^{E_1\oplus N}_{X,Y}(A_1).
\]
Here the algebra in parentheses specifies where the Hall number is computed.
\end{lemma}
\begin{proof}
For a submodule counted on side $i$, the conditions $YJ_i=XJ_i=0$ give
\[
 S\subseteq U\subseteq R\oplus N.
\]
Fix its actual projection $V=\operatorname{pr}_{E_i}U$, so $S\subseteq V\subseteq R$, and put $\Omega=h_0(V)$. Since $\phi$ preserves $\Omega$, also $h_1(V)=\Omega$. Write $q_i=h_i|_V:V\to\Omega$. For $K_0=U\cap N$, put $N'_0=N/K_0$. The submodule $U$ is determined by a unique graph map $g:\Omega\to N'_0$. Its quotient and its own isomorphism type are given by
\begin{equation}\label{eq:graph-labels}
 (T\oplus N'_0)/\{(\omega,g\omega):\omega\in\Omega\}\simeq X,
 \qquad \operatorname{mid}((gq_i)^*\beta_{K_0})\simeq Y,
\end{equation}
where $\beta_{K_0}$ is the class of $0\to K_0\to N\to N'_0\to0$. A minus sign may be placed in the graph and in $g$ simultaneously. For a map $f:W\to N'_0$, the pullback $f^*\beta_{K_0}$ has middle module
$\{(w,n):f(w)=n+K_0\}$; $\operatorname{mid}$ denotes that middle module.

For fixed abstract modules $K,N'$, write $\Ext_C^1(N',K)_N$ for extension classes with middle module isomorphic to $N$. Put $a_W=|\Aut_C(W)|$ in this proof. The contribution of all $K_0$ with these endpoint types is
\begin{equation}\label{eq:framed-graphs}
 \frac{a_N}{a_Ka_{N'}|\Hom_C(N',K)|}
 \#\{(\beta,g):\beta\in\Ext_C^1(N',K)_N,\ g:\Omega\to N',
                  \ \eqref{eq:graph-labels}\text{ holds}\}.
\end{equation}
Indeed, every $K_0$ has $a_Ka_{N'}$ endpoint frames. A fixed extension class has $a_N/|\Hom_C(N',K)|$ exact-sequence representatives on the fixed middle module $N$: the stabilizer is the group of maps $1+iup$, with $u:N'\to K$, where $i,p$ are its two structure maps. Quotient frames transport $g$ at the same time. This proves \eqref{eq:framed-graphs}.

Fix $g$ and choose $\psi$ in \eqref{eq:stable-compensation}. Pullback by $\psi$ is a bijection of $\Ext_C^1(N',K)_N$. Its inverse is pullback by $\psi^{-1}$, and its middle module remains isomorphic to $N$. Moreover,
\[
 (gq_0)^*(\psi^*\beta)=(\psi gq_0)^*\beta
       =(g\phi q_0)^*\beta=(gq_1)^*\beta.
\]
The second equality holds because extensions pull back to zero through projective modules. The quotient condition in \eqref{eq:graph-labels} depends only on $g$. Thus \eqref{eq:framed-graphs} agrees on both sides. Summing over $K,N'$ and over the same submodules $V\subseteq R$ proves the assertion.
\end{proof}

\subsection{Root-isolated pennies}
Let $k$ have characteristic two. The local configurations used below have a root $x$, a loop $r$, and an oriented arm
\[
 x\xrightarrow{a_1}1\xrightarrow{a_2}\cdots
 \xrightarrow{a_{n-1}}n-1\xrightarrow{a_n}x,\qquad n\ge2.
\]
Write $p_j=a_1\cdots a_j$ and $q_i=a_{i+1}\cdots a_n$ for $1\le i,j<n$. The two configurations have parameter $\epsilon\in\{0,1\}$ and relations
\begin{equation}\label{eq:penny}
 a_1\cdots a_n=r^2+\epsilon r^3,\qquad a_na_1=0,
 \qquad r^4=0.
\end{equation}
The root corner has basis $1,r,r^2,r^3$; the root-to-$j$ and $i$-to-root spaces have bases $p_j,rp_j$ and $q_i,q_ir$. In particular,
\[
 r^2p_j=q_ir^2=q_ip_j=0.
\]
The displayed vertices form a full subcategory with exactly the displayed arrows. Nonroot paths through $x$ are generated by $q_irp_j$ and their nonroot multiples. We call the root \emph{isolated} when both corner spaces between $x$ and any outside vertex are zero. Arrows incident with nonroot arm vertices are retained. In particular, $p_jb=0$ for every arrow $b$ from $j$ to an outside vertex. The ray configuration is described in \cite[\S1 and Proposition 2]{Bongartz12}; the nonstandard multiplication is given by \cite[Theorem 10]{Bongartz13}. To obtain \eqref{eq:penny} from the stem coordinates, write its arm arrows as $b_1,\ldots,b_n$. In walking order their relations are
\[
 b_1\cdots b_n=r^2,\qquad b_nb_1=b_nrb_1,\qquad r^4=0.
\]
Multiplication of the second relation by the rest of the arm gives $r^2b_1=r^3b_1$ and, dually, $b_nr^2=b_nr^3$. Since $1+r$ is invertible, $r^2b_1=b_nr^2=0$. Set $a_1=(1+r)b_1$ and $a_i=b_i$ for $2\le i\le n$. Then
\[
 a_1\cdots a_n=r^2+r^3,\qquad a_na_1=0,
 \qquad a_nra_1=b_nrb_1.
\]
Moreover $q_irp_j$ is unchanged, since the additional term contains $r^2b_1=0$. Thus the nonroot relations are preserved. The occurrence of these configurations in representation-finite algebras is used in Section~\ref{sec:global}.

At an isolated split root, $P=e_xA$ is projective-injective with simple socle $S=kr^3$ and top the root simple. Indeed, pair $e_xA$ with $Ae_x$ by the coefficient of $r^3$ in their product. On the root powers the pairing has ones on its antidiagonal and zeros elsewhere. On the bases $p_j,rp_j$ and $q_j,q_jr$ its matrix is
\[
 \begin{pmatrix}\epsilon&1\\1&0\end{pmatrix}.
\]
Root isolation leaves no other blocks. Each block is nondegenerate, and associativity of multiplication makes the pairing a right-module isomorphism $e_xA\simeq D_k(Ae_x)$. Its socle is therefore the root simple, represented by $r^3$. For a Frobenius orbit of such roots, the same assertion follows over the finite ground field: the corresponding projective and injective become isomorphic after a finite splitting extension; restriction of scalars and Krull--Schmidt cancellation give their isomorphism over the ground field.

For several selected roots, their socles form a two-sided ideal $J$. Assume their penny supports are disjoint. The common quotient of the two choices is $C=A/J$. At one split root put $T=e_xC$, $H=\rad T$, and $R=\rad P$. Their spaces are
\[
 T_x=k1\oplus kr\oplus kr^2,\quad H_x=kr\oplus kr^2,
 \quad R_x=kr\oplus kr^2\oplus kr^3,
 \quad T_j=H_j=R_j=kp_j\oplus krp_j.
\]
All outside spaces vanish. Arm arrows propagate the two columns. In $H$, the first arrow sends $r$ to $rp_1$, and the last sends $p_{n-1}$ to $r^2$ and $rp_{n-1}$ to zero. The loop sends $r$ to $r^2$.

Identify the two radicals by the identity on arms and by
\begin{equation}\label{eq:radical-identification}
 r\longmapsto r+r^2,\qquad r^2\longmapsto r^2+r^3,
 \qquad r^3\longmapsto r^3
\end{equation}
at the root, from the nonstandard to the standard side. The last-arrow value $r^2+r^3$ is thereby sent to $r^2$; the loop and first arrow also commute. Identify $P/S$ by the path bases. The two boundary maps $R\to T$ differ by $\phi$, which acts on $H_x$ as $v\mapsto v+vr$ and is the identity elsewhere. This is a $C$-automorphism preserving every submodule of every $H^m$, for $m\ge0$.

\begin{lemma}\label{lem:isolated-compensation}
For $\Omega\subseteq H^m$ and $g:\Omega\to N'$, condition \eqref{eq:stable-compensation} holds with $\psi=1+n$, where $n^2=0$, $n$ kills $\rad N'$, and its image lies in the root part of $\soc N'\cap\rad N'$.
\end{lemma}
\begin{proof}
Write a root vector as $r\xi+r^2\eta$, with $\xi,\eta\in k^m$. Let $W\subseteq k^m$ be the space of first coefficients of vectors in $\Omega_x$. The first-coefficient map has kernel $L=\ker(r:\Omega_x\to\Omega_x)$. Choose a linear section $\sigma:W\to\Omega_x$. Closure under paths gives $r^2W\subseteq\Omega_x$ and $rp_jW\subseteq\Omega_j$.

On the arm chain $1\to\cdots\to n-1$, the second columns $rp_jW$ form a constant subrepresentation. It is injective in the category of chain representations: for one copy, a map from a chain $V$ is determined by an arbitrary functional on $V_{n-1}$, and the earlier components are its composites with the intervening arrows. This Hom functor is exact. Hence there are chain-compatible retractions $\rho_j:\Omega_j\to W$ with $\rho_j(rp_j\xi)=\xi$.

We first construct a projective factorization. Suppose $u:W\to(\rad N')_x$ satisfies $u(\xi)r^2=u(\xi)rp_j=0$. The map with root values $\sigma(\xi)\mapsto u(\xi)r$, zero on $L$ and zero on the arms, factors through projectives. Indeed, root isolation gives
\[
 (\rad N')_x=N'_xr+\sum_{i=1}^{n-1}N'_iq_i.
\]
Choose linear $w:W\to N'_x$, $z_i:W\to N'_i$ with $u(\xi)=w(\xi)r+\sum_i z_i(\xi)q_i$. The $w(\xi)r^2$ term factors through $T\otimes_kW$: send $\sigma(\xi)$ to $r^2\otimes \xi$, all other components to zero, and then $1\otimes\xi$ to $w(\xi)$. This commutes with the first arrow because $r^2p_1=0$ and with the last because its source image lies in $L$.

For the $i$th remaining term use $e_iC\otimes_kW$. Send $\sigma(\xi)$ to $q_ir\otimes \xi$, send $L$ to zero, and send $v\in\Omega_j$ to
\[
 q_irp_j\otimes\rho_j(v).
\]
The loop commutes by $q_ir^2=0$; the first arrow commutes by the retraction identity; internal arrows commute by chain compatibility. The last arrow gives $q_irp_{n-1}a_n=q_ir^3=0$, and its source image lies in $L$. Outside arrows commute because $p_jb=0$ and the source has zero outside spaces. Send the projective generator $e_i\otimes\xi$ to $z_i(\xi)$. All these factors use the same retraction, so their total arm value is zero:
\[
 \sum_i z_i(\eta)q_irp_j=u(\eta)rp_j-w(\eta)r^2p_j=0\qquad(\eta\in W).
\]
Their root sum is $u(\xi)r$, as required.

Now let $\pi:N'_x\to N'_x/(\rad N')_x$ be the root top projection. Put $B=\pi(gL)$ and define $t:W\to N'_x/(\rad N')_x/B$ by $t(\xi)=\pi(g\sigma(\xi))+B$. Choose $W=W_0\oplus W_v$ with $W_0=\ker t$, and denote the projections by $p_0,p_v$. The spaces $B$ and $\pi(g\sigma(W_v))$ are disjoint. Define a map from the root top to the root socle by zero on $B$, by
\[
 \pi(g\sigma(\xi))\longmapsto g\sigma(\xi)r\qquad(\xi\in W_v),
\]
and by zero on a complementary space. The values lie in $\rad N'\cap\soc N'$ because $\sigma(\xi)r^2=\sigma(\xi)rp_1=0$ and the root has no other outgoing arrow. Composing with the top projection and socle inclusion gives a $C$-endomorphism $n$ with $n^2=0$. It satisfies $ngL=0$ and $ng\sigma(\xi)=g\sigma(p_v\xi)r$.

For $\xi\in W_0$, choose linearly $\ell(\xi)\in L$ with $\pi(g\sigma(\xi))=\pi(g\ell(\xi))$. Then $u(\xi)=g(\sigma(\xi)-\ell(\xi))$ lies in $(\rad N')_x$, satisfies $u(\xi)r=g\sigma(\xi)r$, and is killed by $r^2$ and $rp_j$. Extend it by $u p_0$. The factorization just constructed gives
\[
 g(\phi-1)-ng\in\mathcal P_C(\Omega,N').
\]
Thus $\psi=1+n$ has the asserted properties.
\end{proof}

\begin{lemma}\label{lem:compensation-descent}
Suppose the common boundary data and their shear are defined over a finite field $k$, and become the preceding disjoint root data after scalar extension. Then \eqref{eq:stable-compensation} holds over $k$ and over every finite extension of $k$.
\end{lemma}
\begin{proof}
Let $\topm_{\rm sel}N'$ and $\soc_{\rm sel}N'$ be the sums of the selected simple types in the top and socle. Elements of
\[
 \Hom_C(\topm_{\rm sel}N',\rad N'\cap\soc_{\rm sel}N')
\]
give square-zero endomorphisms $n$ by projection and inclusion. The condition $ng-g(\phi-1)\in\mathcal P_C(\Omega,N')$ is linear in $n$. Moreover,
\[
 \mathcal P_C(\Omega,N')=\im\bigl(
 \Hom_C(\Omega,C)\otimes_k\Hom_C(C,N')\longrightarrow\Hom_C(\Omega,N')\bigr),
\]
where the map is composition: a projective is a summand of a finite free module. This image commutes with field extension, as do Hom spaces, radicals and the indicated semisimple parts over a perfect field. Over a splitting field, $\Omega$ decomposes by the disjoint penny supports, since arrows between these supports act by zero in its ambient module. Apply Lemma~\ref{lem:isolated-compensation} to each summand and add the corrections. The resulting linear system has a solution after extension, hence over $k$. Then $1+n$ is the required automorphism. The argument applies over any finite extension as the new base field.
\end{proof}

\begin{proposition}\label{prop:isolated-comparison}
Suppose $A,A^0$ have the actual common quotient and root-isolated boundary data above. Retain every $C$-module and replace each selected projective-injective $P_s$ by its corresponding $P_s^0$, with the same multiplicity. This correspondence preserves all Hall numbers over every finite extension. If $A^0$ is geometrically standard and representation-finite, then $A$ has Hall polynomials.
\end{proposition}
\begin{proof}
For an indecomposable $A$-module $M$ with $MJ\ne0$, a map $P_s\to M$ is nonzero on the simple essential socle of some $P_s$. Its kernel is zero. Injectivity of $P_s$ splits the embedding, so $M\simeq P_s$. All other indecomposables are exactly the $C$-modules. The same holds on the other side. Lemmas~\ref{lem:graph-comparison} and \ref{lem:compensation-descent} give the Hall comparison when both endpoints are $C$-modules and the middle module is arbitrary.

Hom dimensions agree under the correspondence. For a $C$-module $V$,
\[
 \Hom_A(P_s,V)\simeq Ve_s,\qquad
 \Hom_A(V,P_s)\simeq D_k(Ve_s).
\]
The second identity uses $P_s\simeq D_k(Ae_s)$, since it is the injective envelope of the root simple. The dual of a finite residue field is isomorphic to itself as a module over that field. Hom spaces between selected projectives are the corner spaces; these have the dimensions of the common quotient corners plus their corresponding socle corners. Residue-field degrees agree for $C$-labels, and those of $P_s,P_s^0$ equal that of their common top. Formula~\eqref{eq:aut} therefore gives equal automorphism orders.

Write $X=E_X\oplus X_c$, $Y=E_Y\oplus Y_c$ with $X_c,Y_c$ annihilated by $J$ and $E_X,E_Y$ sums of selected projective-injectives. Projectivity and injectivity give
\[
 \Ext_A^1(X,Y)\simeq\Ext_A^1(X_c,Y_c),
\]
with middle module $E_X\oplus E_Y$ plus the middle module on the right. Cancel these summands in a prescribed middle module; if their multiplicities are insufficient, the extension set is empty. For arbitrary endpoints the exact-sequence count is
\begin{equation}\label{eq:extension-hall}
 F^M_{X,Y}=
 \frac{|\Ext_A^1(X,Y)_M|\,|\Aut_A(M)|}
 {|\Aut_A(X)|\,|\Aut_A(Y)|\,|\Hom_A(X,Y)|}.
\end{equation}
It follows by the same endpoint-frame count as \eqref{eq:framed-graphs}. Apply it first to the $C$-endpoints and then to $X,Y$ to obtain the comparison. Over an arbitrary finite extension repeat the argument with its indecomposable decomposition.

The indecomposable correspondence preserves all residue-field degrees, so $D_A=D_{A^0}$. Theorem~\ref{thm:standard} supplies the polynomial on $A^0$ over exactly the conservative extensions required for $A$.
\end{proof}
